\documentclass[a4paper, 12pt]{article}

%Class
	%\documentclass[a4paper,12pt]{article}

%Packages
	\usepackage{amsmath}					% Standard maths
	\usepackage{amssymb}					% Standard maths
	\usepackage{amsthm}						% Standard maths
	\usepackage{thmtools, thm-restate}% Theorem styles
	\usepackage{mathtools}
	\usepackage[newzealand]{babel}% Date/time in British format (yes, I know it says NZ)
	\usepackage{datetime}					% Date/time
	\usepackage{multirow}					% For stacked subscripts
	\usepackage{xcolor}						% For colour!
	\usepackage[normalem]{ulem}		% For strikeout
	\usepackage{isomath}
	\usepackage{fancyhdr}					% For the headers
	\usepackage{mfirstuc}					% For the headers
	\usepackage[hmargin=1in,top=0.8in,bottom=1in]{geometry} % Page margins more reasonable
	\usepackage{graphicx}					% For including graphics
	%\usepackage{float}						% For something
	\usepackage{enumitem}					% For numbering tweaks
	\usepackage[margin=15pt,font={footnotesize},labelfont=bf,format=hang]{caption}
\usepackage[margin=15pt,font={footnotesize},labelfont=bf,format=hang]{subcaption}
	\usepackage{url}
	\usepackage[hidelinks]{hyperref}	% For hyperlinks
	\usepackage{breakurl}
	\usepackage{cleveref}
	\usepackage[scaled=.92]{helvet}
	\newbool{show_question}
	\newbool{show_solution}
	\newcommand{\solutioninheading}{\ifbool{show_solution}{: SOLUTIONS}{}}
	\usepackage{pgfplots}
\newcommand{\ep}{\varepsilon}

\makeatletter 
\newcommand\mynobreakpar{\par\nobreak\@afterheading} 
\makeatother

    \newcommand{\question}[3]{
    	\item \textbf{#1} \mynobreakpar \nobreak \medskip % Title
    	%
    	\ifbool{show_question}{\nobreak#2}{} % Question
    	%
    	%
    	\ifbool{show_question}{\ifbool{show_solution}{\par\medskip\textbf{Solution:}}{}}{} % Both
    	\ifbool{show_solution}{#3}{} % Answer
    }		

%MATHEMATICAL SHORTCUTS
%Two and three-part definitions
	\newcommand{\twopartdef}[4]
	{	\left\{
			\begin{array}{ll}
				#1 & \mbox{if } #2 \\
				#3 & \mbox{if } #4
			\end{array}
		\right.}
%	\newcommand{\threepartdef}[6]
%	{	\left\{
%			\begin{array}{ll}
%				#1 & \mbox{if } #2 \\
%				#3 & \mbox{if } #4 \\
%				#5 & \mbox{if } #6
%			\end{array}
%		\right.}
%Blackboard bold	
	\newcommand{\N}{{\mathbb N}} 
	\newcommand{\R}{{\mathbb R}} 
	\newcommand{\C}{{\mathbb C}} 
	%Curly letters
%	\newcommand{\E}{{\mathcal E}}	
%	\newcommand{\F}{{\mathcal F}} 
%	\newcommand{\Null}{{\mathcal N}} 
%	\newcommand{\R}{\mathcal{R}} 
%	\newcommand{\Z}{\mathcal{Z}} 
%Uprights
	\newcommand{\D}{{\mathrm D}}
	\renewcommand{\d}{{\mathrm d}}
%Easy Greeks
%	\def\a{\alpha}
%	\def\g{\gamma}
%	\newcommand{\e}{{\varepsilon}}
%	\newcommand{\la}{{\lambda}}  
%	\newcommand{\om}{{\Omega}}
	\renewcommand{\epsilon}{\varepsilon}
	
%Vector operators
	\def \p {\partial}
	\newcommand{\del}{{\nabla}}
	\newcommand{\grad}{{\nabla}}
    \newcommand{\vdel}{\boldsymbol{\nabla}}
    \newcommand{\vgrad}{\boldsymbol{\nabla}}
    \newcommand{\vnabla}{\boldsymbol{\nabla}}
	\newcommand{\cross}{{\times}}
	\newcommand{\Lap}{{\nabla^2}} 
%Arrows
%	\newcommand{\goesto}[1]{\xrightarrow[#1]{}}
%hats
%	\newcommand{\nhat}{\widehat{\v{n}}}
%	\newcommand{\shat}{\widehat{\v{s}}}
	\renewcommand{\hat}{\widehat}
%Note to self
%	\newcommand{\note}[1]{[\textbf{\textsl{\MakeUppercase{#1}}}]}
	
%Stress tensor
%	\newcommand{\T}{\vg{\sigma}}

% Vectors, quaternions etc.
\renewcommand{\v}[1]{\mathbfit{#1}}      % Generic vector
\newcommand{\vc}[1]{\bm{\mathcal{#1}}}      % vector mathcal
\newcommand{\uv}[1]{\widehat{\mathbfit{#1}}}      % Generic unit vector
\newcommand{\q}[1]{\mathsfbfit{#1}}        % Generic quaternion
\renewcommand{\t}[1]{\mathsfbfit{#1}}	 % Generic tensor
\newcommand{\m}[1]{\mathsfbfit{#1}}	 % Generic matrix
\newcommand{\vu}[1]{\mathbf{#1}}         % Upright vector (for zero vector)
\newcommand{\qu}[1]{\mathbf{#1}}       % Upright quaternion (for zero quaternion)
\newcommand{\tu}[1]{\bm{\mathsf{#1}}}	 % Upright tensor (for zero tensor)

%Real and imaginary
\renewcommand{\Re}{\operatorname{Re}}
\renewcommand{\Im}{\operatorname{Im}}
	
%Non-dim numbers
%	\renewcommand{\Re}{\operatorname{Re}}
%	\newcommand{\Bi}{\operatorname{Bi}}
%	\newcommand{\Fo}{\operatorname{Fo}}
	
%Misc
\DeclareMathSymbol{\Delta}{\mathalpha}{operators}{1} % Keeps Delta upright
\DeclareMathSymbol{\Gamma}{\mathalpha}{operators}{0} % Keeps Gamma upright
	\newcommand{\cosec}{\operatorname{cosec}}
	\newcommand{\cosech}{\operatorname{cosech}}
	\newcommand{\sech}{\operatorname{sech}}
	\newcommand{\tr}{\operatorname{tr}}
	\newcommand{\erf}{\operatorname{erf}}
	\newcommand{\order}{\mathcal{O}}
%	\newcommand{\V}{\mathcal{V}} % 	CurlyV = (Bi V) / (1 + Bi)

% Differentiating 
\newcommand{\fd}[2]{\mathchoice{\frac{\d #1}{\d #2}}{\d #1/\d #2}{\d #1/\d #2}{\d #1/\d #2}}
\newcommand{\fdd}[2]{\mathchoice{\frac{\d^2 #1}{\d #2^2}}{\d^2 #1/\d #2^2}{\d^2 #1/\d #2^2}{\d^2 #1/\d #2^2}}
\newcommand{\fddd}[2]{\mathchoice{\frac{\d^3 #1}{\d #2^3}}{\d^3 #1/\d #2^3}{\d^3 #1/\d #2^3}{\d^3 #1/\d #2^3}}
\newcommand{\fdddd}[2]{\mathchoice{\frac{\d^4 #1}{\d #2^4}}{\d^4 #1/\d #2^4}{\d^4 #1/\d #2^4}{\d^4 #1/\d #2^4}}
\newcommand{\fdn}[3]{\mathchoice{\frac{\d^{#1} #2}{\d #3^{#1}}}{\d^{#1} #2/\d #3^{#1}}{\d^{#1} #2/\d #3^{#1}}{\d^{#1} #2/\d #3^{#1}}}
\newcommand{\pd}[2]{\mathchoice{\frac{\p #1}{\p #2}}{\p #1/\p #2}{\p #1/\p #2}{\p #1/\p #2}}
\newcommand{\pdd}[2]{\mathchoice{\frac{\p^2 #1}{\p #2^2}}{\p^2 #1/\p #2^2}{\p^2 #1/\p #2^2}{\p^2 #1/\p #2^2}}
\newcommand{\pddmixed}[3]{\frac{\p^2 #1}{\p #2 \p #3}}
\newcommand{\pddd}[2]{\frac{\p^3 #1}{\p #2^3}}
\newcommand{\pdn}[3]{\mathchoice{\frac{\p^{#1} #2}{\p #3^{#1}}}{\p^{#1} #2/\p #3^{#1}}{\p^{#1} #2/\p #3^{#1}}{\p^{#1} #2/\p #3^{#1}}}
	
	% Misc
\newcommand{\e}{\mathrm{e}}
\renewcommand{\i}{\mathrm{i}}
\newcommand{\dx}{\,\mathrm{d}x}
\newcommand{\dy}{\,\mathrm{d}y}
\renewcommand{\epsilon}{\varepsilon}
\renewcommand{\Re}{\operatorname{Re}}
\newcommand{\Four}{\mathcal{F}}

\newcommand{\mat}[4]{\begin{pmatrix}#1 & #2 \\ #3 & #4 \end{pmatrix}}

	%Column vectors
	\makeatletter
\newcommand\rcvector[2][\\]{\ensuremath{%
  \global\def\rc@delim{#1}%
    \negthinspace\begin{pmatrix}
      \rc@vector #2,\relax\noexpand\@eolst%
    \end{pmatrix}}}
\def\rc@vector #1,#2\@eolst{%
  \ifx\relax#2\relax
    #1
  \else
    #1\rc@delim
    \rc@vector #2\@eolst%
  \fi}
\makeatother

\newcommand{\colvec}{\rcvector}
\newcommand{\rowvec}{\rcvector}
\renewcommand{\binom}{\rcvector}
	
	%Colours
%	\newcommand{\orange}[1]{{\color{orange} #1 }}
%	\newcommand{\blue}[1]{{\color{blue} #1 }}

%Theorem styles
%	\declaretheoremstyle[headpunct={\:\:\:}, shaded={bgcolor={rgb}{0.9,0.9,0.9}, margin=5pt}]{problem}
%	\declaretheoremstyle[headpunct={\:\:\:}, shaded={rulecolor=black, rulewidth=0.4pt, bgcolor={rgb}{1,1,1}, margin=5pt}]{definition}
%	\declaretheoremstyle[headpunct={\:\:\:}, spaceabove=15pt, notefont=\bf, notebraces={(}{)}]{theorem}
%	\declaretheoremstyle[headpunct={\:\:\:}, numbered=no, spaceabove=15pt, qed={\checkmark}]{solution}
%	
%	\declaretheorem[style=definition,numberwithin=section]{definition}
%	\declaretheorem[numberlike=definition, style=theorem]{theorem}
%	\declaretheorem[numberlike=definition, style=theorem]{lemma}
%	\declaretheorem[numberlike=definition,style=problem]{problem}
%	\declaretheorem[numberlike=definition, style=theorem]{proposition}
%	\declaretheorem[numberlike=definition, style=theorem]{corollary}
%	\declaretheorem[numberlike=definition, style=theorem]{remark}
%	\declaretheorem[numberlike=definition, style=theorem]{example}
%	\declaretheorem[numberlike=definition, style=theorem]{notation}
%	\declaretheorem[style=solution]{solution}

%Depth of display breaks to allow	
	\allowdisplaybreaks[1]

%How deep do we want the Table of Contents to go?	
%	\setcounter{tocdepth}{1}

%Sections
\makeatletter
\renewcommand{\section}{\@startsection
{section}%                   % the name
{1}%                         % the level
{0mm}%                       % the indent
{-\baselineskip}%            % the before skip
{0.5\baselineskip}%          % the after skip
{\normalfont\bfseries}} % the style
\makeatother

% Wavy underline
\makeatletter
\font\sixly=lasyb10 scaled 1200
\def\tinners{\bgroup \markoverwith{\lower8.5\p@\hbox{\sixly \char58}}\ULon}
\makeatother
\newcommand{\answer}[1]{\tinners{\; #1 \;}}

%Setting the footnotes to use *,**,+ etc rather than 1,2,3	
\renewcommand{\thefootnote}{\fnsymbol{footnote}}

%Italic quote environment
%\newenvironment{quoteitalic}{\begin{quote}\itshape}{\end{quote}}

%Heading
\newcommand{\heading}[2]{
\setlength{\parindent}{0pt}
\textbf{MATH3171 (Mathematical Biology III)}

\textbf{YEAR 2025--26, \MakeUppercase{#1} TERM}
\setlength{\parskip}{3ex}

\textbf{\MakeUppercase{#2}}
%
\setlength{\parskip}{2ex}

}

%========================================================================================================================
%DOCUMENT BEGINS HERE!

\setbool{show_question}{true}
\setbool{show_solution}{false}
\newcommand{\mvec}[2]{\left(
    \begin{array}{c}#1\\#2
\end{array}\right)}
\begin{document}
%
\heading{Epiphany}{Additional problem sheet 4\solutioninheading}
\ifbool{show_question}{\textbf{Topics covered}: Chapters 10--11.}{}

\begin{enumerate}[label={\bf \arabic{*}.\;}, ref={\arabic{*}},itemsep=3ex, left=0ex]

    \question{Turing conditions (Exam section A type)}
    {
      Consider the following \emph{non-dimensionalised} reaction--diffusion equation for the scalar densities $u$ and $v$:
      \begin{subequations}
        \label{rds2}
        \begin{align}
          \pd{u}{t} &= \nabla^2 u + \gamma F(u,v),\\
          \pd{v}{t} &= D \nabla^2 v + \gamma G(u,v),
        \end{align}
      \end{subequations}
      where $D>0$ and $\gamma>0$ are constants. We assume \emph{no-flux} boundary conditions on a Cartesian domain $[0,L_1]\times[0,L_2]$ with coordinates $(x_1,x_2)$. The Turing conditions for pattern formation are
      \begin{align}
        & F_u +G_v <0, \quad F_uG_v  -G_uF_v> 0,\nonumber \\
        & G_v + DF_u>0,\quad (G_v+ D F_u)^2 - 4 D (F_uG_v  -G_uF_v) >0,\label{turingconditions2}
      \end{align}
      where we have used the the notation
      \begin{align*}
        &F_{u} = \pd{F}{u}\bigg\vert_{u=u_0,v=v_0},\quad F_{v} = \pd{F}{v}\bigg\vert_{u=u_0,v=v_0},\\
        &G_{u} = \pd{G}{u}\bigg\vert_{u=u_0,v=v_0},\quad G_{v} = \pd{G}{v}\bigg\vert_{u=u_0,v=v_0}.
      \end{align*}
      \begin{enumerate}[label=(\alph*)]
        \item{State the conditions for a homogeneous equilibrium of \cref{rds2}.
          }
        \item{  Define what is meant by a pattern in this context, give the explicit mathematical form for the pattern and describe what the conditions in \cref{turingconditions2} enforce.
          }
        \item{State a condition on $D$ under the assumption that $F_u>0.$
          }
      \end{enumerate}
    }
    {
      \begin{enumerate}[label=(\alph*)]
        \item The homogeneous equilibrium satisfies
          \begin{equation}
            F(u,v) =0,\quad G(u,v) = 0.
          \end{equation}

        \item This homogeneous equilibrium is unstable to small variations in density. This instability develops into an inhomogeneous pattern in the form often of spots or stripes (this being a Cartesian grid). The linearised version of the main equations has solutions in the form
          \begin{equation}
            \sum_{i=0}^{\infty}\sum_{j=0}^{\infty}\e^{\lambda(k_s^{ij}) t}\cos\left(\frac{ i \pi}{L_1}x_1\right)\cos\left(\frac{ j\pi}{L_2}x_2\right),
          \end{equation}
          where
          \begin{equation}
            k_s^{ij} = \frac{ i^2\pi^2}{L_1^2}+\frac{j^2\pi^2}{L_2^2}
          \end{equation}
          is the pattern number.
          The function $\lambda(k_s)$ dictates if the mode decays exponentially $\lambda(k_s)<0,$ or grows, $\lambda(k_s)\geq 0$. The growing modes form the pattern.

          The first two conditions,
          \begin{equation}
            F_u +G_v <0, \quad F_uG_v  -G_u F_v> 0,
          \end{equation}
          ensure that the $k_s^{00}=0$ mode, the homogeneous perturbation, always decays.

          The other two conditions,
          \begin{equation}
            G_v + DF_u>0,\quad (G_v+ D F_u)^2 - 4 D (F_uG_v  -G_uF_v)>0,
          \end{equation}
          ensure there is some finite set of wave modes $k_s\in[k_s^{\min},k_s^{\max}]$ for which $\lambda(k_s^{ij})>0$ and hence a pattern can form.

        \item
          If $F_u>0$ the first condition requires $G_v<0$ and $\vert G_v\vert> \vert F_u\vert$. The third condition then shows
          \begin{equation}
            D \vert F_u\vert - \vert G_v\vert >0 \implies
            D>  \frac{\vert G_v\vert}{\vert F_u\vert},
          \end{equation}
          and hence $D>1$.

      \end{enumerate}

    }
    %

    %===========================================================================

    %
    \question{Periodic boundary conditions (Exam section B type)}
    {
      Consider the following generic reaction--diffusion system,
      \begin{subequations}
        \begin{align}
          \pd{u}{t} &= \nabla^2 u + \gamma F(u,v),\\
          \pd{v}{t} &=  D\nabla^2 v + \gamma G(u,v),
        \end{align}
      \end{subequations}
      with $D$ and $\gamma$ positive constants. We assume a Cartesian domain $[0,L_1]\times[0,L_2]$ with coordinates $(x_1,x_2)$. We assume the densities $u$ and $v$ have periodic boundary conditions
      \begin{align}
        u(0,x_2) &= u(L_1,x_2), & u(x_1,0) &= u(x_1,L_2),\nonumber\\
        v(0,x_2) &= v(L_1,x_2), & v(x_1,0) &= v(x_1,L_2).
      \end{align}
      \begin{enumerate}[label=(\alph*)]
        \item{State what is meant by a homogeneous equilibrium and state the equations of such equilibria in this case.
          }
        \item{
            Consider an expansion $u(x,t)= u_0 + \epsilon u_1(x,t)$ and $v(x,t)= v_0 + \epsilon v_1(x,t)$, show that, to $\order(\epsilon)$, solutions to $u_1$ and $v_1$ take the form
            \begin{align}
              \sum_{i=0}^{\infty}\sum_{j=0}^{\infty}\e^{\lambda(i,j)t}\left[C_1^{ij}\sin\left(\frac{i\pi}{L_1}x_1\right)\sin\left(\frac{j\pi}{L_2}x_2\right) + C_2^{ij}\sin\left(\frac{i\pi}{L_1}x_1\right)\cos\left(\frac{j 2\pi}{L_2}x_2\right)\right.\nonumber \\
              \left. + C_3^{ij}\cos\left(\frac{i 2 \pi}{L_1}x_1\right)\sin\left(\frac{j \pi}{L_2}x_2\right)  + C_4^{ij}\cos\left(\frac{i 2\pi}{ L_1}x_1\right)\cos\left(\frac{j 2\pi}{L_2}x_2\right) \right],
            \end{align}
            where the $C_\alpha^{ij}$ are constants. You must state the explicit form of the functions $\lambda(i,j)$.
          }
        \item{What would the linearised solution look like if we had instead chosen Dirichlet boundary conditions, i.e.\
            \begin{align}
              u(x_1,0) &= u(x_1,L_2)  = u(0,x_2) = u(L_1,x_2)= u_0,\nonumber\\
              v(x_1,0) &= v(x_1,L_2)  =v(0,x_2) = v(L_1,x_2)= v_0,
            \end{align}
            where $u_0$ and $v_0$ are the homogeneous equilibrium values?
          }
        \item{We have now considered two distinct sets of boundary conditions. In which case are the Turing conditions
            \begin{equation}
              F_u  + G_v<0, \quad F_uG_v  + F_v G_u>0,
            \end{equation}
            necessary for  pattern formation? Justify your answer with respect to the linear solutions for $u_1$ in both cases.
          }
      \end{enumerate}
    }
    {
      \begin{enumerate}[label=(\alph*)]
        \item{
            A homogeneous equilibrium is one for which both spatial and temporal derivatives vanish. In the case at hand the equations of a homogenous equilibrium $(u_0,v_0)$ are
            \begin{equation}
              F(u_0,v_0)  =G(u_0,v_0) = 0.
            \end{equation}
          }
        \item{
            To $\order(\epsilon)$ the main equations are
            \begin{align}
              \pd{u_1}{t} &= \nabla^2 u_1 + \gamma\left(F_u u_1 + F_v v_1\right),\\
              \pd{v_1}{t} &= D\nabla^2 v_1 + \gamma\left(G_u u_1 + G_v v_1\right),
            \end{align}
            where, for example, $F_u = \pd{F}{u}$, evaluated at $u=u_0$.
            We seek solutions in the form $u_1 = \Re\left(A\e^{\i\v{k}\cdot  x + \lambda t}\right)$ and $v_1 = \Re\left(B\e^{\i\v{k}\cdot  x + \lambda t}\right)$; then our linearised equations take the form
            \begin{equation}
              \label{lineq}
              \lambda
              \begin{pmatrix}u\\v
              \end{pmatrix} = \m{J}
              \begin{pmatrix}u \\ v
              \end{pmatrix},
            \end{equation}
            where
            \begin{equation}
              \m{J} =
              \begin{pmatrix}
                -(k_1^2 + k_2^2) + \gamma F_u & \gamma F_v \\
                \gamma G_u & -D(k_1^2 + k_2^2) + \gamma G_v
              \end{pmatrix}.
            \end{equation}
            We now apply the boundary conditions. Consider the $x_1$ boundary conditions. Depending on the choice of the constants $A$ (or $B$) the solutions will take the form
            \begin{equation}
              A\e^{\i k_2 x_2 + \lambda t}\sin(k_1x_1) \quad \mbox{ or } \quad A\e^{\i k_2 x_2 + \lambda t}\cos(k_1x_1).
            \end{equation}
            For the sine case the boundary conditions require
            \begin{equation}
              \sin(0) = \sin(k_1 L_1),
            \end{equation}
            so $k_1 L_1 = n \pi$ for $n\in 0,1,2,3\dots$.

            For the cosine case we require
            \begin{equation}
              \cos(0) = \cos(k_1 L_1),
            \end{equation}
            so $k_1 L_1 =  n 2\pi$  for $n\in 0,1,2,3\dots$.

            One can easily apply this argument to the $x_2$ direction. The quoted solution is then the sum over all permissible combinations.

            Finally, in order to solve \cref{lineq} for nontrivial solutions $u_1$ and $v_1$ we require $\det(\m{J}-\lambda \m{I}=0$ where $\m{I}$ is the identity matrix. Solutions to this take the form
              \begin{equation}
                \lambda(i,j) = \frac{1}{2}\left[\tr\m{J} \pm \sqrt{\left(\tr\m{J}\right)^2 - 4\det\m{J}}\right]
              \end{equation}
              with
              \begin{align}
                \tr\m{J} &= \gamma(F_u + G_v)  - k_s(i,j)(D+1),\\
                \det\m{J} & =D k_s(i,j)^2 - \gamma k_s(i,j) \left(D F_u + G_v\right)  + \gamma^2(F_uG_v- F_vG_u),
              \end{align}
              and
              \begin{align}
                k_s(i,j) = \pi^2\left( \frac{i^2}{L_1^2}+ \frac{j^2}{L_2^2}\right),\\
                k_s(i,j) = \pi^2\left( \frac{i^2}{L_1^2}+ \frac{4 j^2}{L_2^2}\right),\\
                k_s(i,j) = \pi^2\left( \frac{4i^2}{L_1^2}+ \frac{j^2}{L_2^2}\right),\\
                k_s(i,j) = \pi^2\left( \frac{4i^2}{L_1^2}+ \frac{4 j^2}{L_2^2}\right),
              \end{align}
              depending on the particular solution type.
            }
          \item{
              In this case the linearised solutions $u_1$ and $v_1$ must satisfy
              \begin{align}
                u_1(x_1,0) = u_1(x_1,L_2)  = u_1(0,x_2) = u_1(L_1,x_2) &= 0,\\   v_1(x_1,0) = v_1(x_1,L_2)  =v_1(0,x_2) = v_1(L_1,x_2) &= 0.
              \end{align}
              In this case the cosine form of the solutions cannot satisfy the boundary conditions as $\cos(0)=1$. The sine boundary condition then require
              \begin{equation}
                k_1 = \frac{n\pi}{L_1} \quad \text{and} \quad k_2 = \frac{m\pi}{L_2},
              \end{equation}
              so that
              \begin{align}
                \sum_{i=1}^{\infty}\sum_{j=1}^{\infty}\e^{\lambda(i,j)t}C_1^{ij}\sin\left(\frac{i\pi}{L_1}x_1\right)\sin\left(\frac{j\pi}{L_2}x_2\right).
              \end{align}
              The summation accounts for the fact that the zeroth mode $i=j=0$ is $0$. The form of $\lambda(i,j)$ will be as in the periodic case with
              \begin{align}
                k_s(i,j) = \pi^2\left( \frac{i^2}{L_1^2}+ \frac{j^2}{L_2^2}\right).
              \end{align}
            }
          \item{
              These two conditions enforce that constant changes in the system, the $i=j=0$ modes, do not grow in time. This is necessary to ensure any potential patterns are not drowned out. In the Dirichlet case (form the previous part) there is no non-zero $i=j=0$ mode, so we have no need to enforce the conditions. For the periodic case the $\cos()\cos()$ solution is potentially non-zero (if $C_4^{00}\neq 0$) so we must enforce these two conditions.
            }
        \end{enumerate}

      }
      %

      %===========================================================================

      %
      \question{Nonlinear diffusion (Exam section B type)}
      {
        Consider the following variant of a reaction--diffusion system,
        \begin{subequations}
          \begin{align}
            \pd{u}{t} &= \nabla^2 u +  F(u,v) + \vnabla u \cdot \vnabla (v^2),\\
            \pd{v}{t} &=  D\nabla^2 v +  G(u,v)+ v\nabla^2 u,
          \end{align}
        \end{subequations}
        with $D>0$ constant. We assume a Cartesian domain $[0,L_1]\times[0,L_2]$ with coordinates $(x_1,x_2)$, and we also assume the densities $u$ and $v$ satisfy no-flux boundary conditions.
        \begin{enumerate}[label=(\alph*)]
          \item{State the equations of homogeneous equilibria of this system
            }
          \item{
              State explicitly the boundary conditions of the system.
            }
          \item{Show that the Turing conditions required for pattern formation take the form
              \begin{align}
                &F_u + G_v<0,\quad  F_uG_v- F_vG_u>0,\nonumber\\
                &D F_u + G_v - v_0 F_v >0, \quad \left(D F_u + G_v - v_0 F_v\right)^2 - 4D(F_uG_v- F_vG_u)>0.
              \end{align}
            }
          \item{
              Consider the case
              \begin{equation}
                F(u,v) = u(1-u) - uv,\quad G(u,v)= v(1-v) + \gamma u v,
              \end{equation}
              where $\gamma$ is a real constant. Find all the equilibria for which at least one of  $u$ or $v$ is $0$.
            }
          \item{For each of the equilibria found in the previous question, establish whether the Turing conditions can be satisfied. If they can be satisfied state the  domain of $\gamma$ values for which patterns can form.
            }
          \item{\emph{Good practice bonus question:} Find the other possible nonzero equilibrium and state conditions on $\gamma$ (if any) for which permit pattern formation. I am happy to look at/discuss your solutions if you bring them to me.}
        \end{enumerate}
      }
      {
        \begin{enumerate}[label=(\alph*)]
          \item Homogeneous equilibria satisfy
            \begin{equation}
              F(u,v) = 0,\quad G(u,v)=0.
            \end{equation}

          \item The boundary conditions are
            \begin{align}
              \pd{u}{x_1}(0,x_2)= \pd{u}{x_1}(L_1,x_2)=  \pd{u}{x_2}(x_1,0) = \pd{u}{x_2}(x_1,L_2)=0,\\
              \pd{v}{x_1}(0,x_2)= \pd{v}{x_1}(L_1,x_2)=  \pd{v}{x_2}(x_1,0) = \pd{v}{x_2}(x_1,L_2)=0.
            \end{align}
            I would accept other notation as long as it is clear, e.g.
            \begin{equation}
              \pd{u}{x_1}\bigg\vert_{x_1=0} = 0,\; \forall x_2\in[0,L_2].
            \end{equation}

          \item{
              We assume the homogeneous equilibria are labelled $u_0$  and $v_0$. We assume $u(x,t) = u_0+\epsilon u_1(x,t)$ and $v(x,t) = v_0+\epsilon v_1(x,t)$. We expand the main equations; the linear and $F$, $G$ expansions are clear and obvious. For the nonlinear terms we have
              \begin{align}
                \vnabla (u^2) &= \vnabla(u_0 + \epsilon 2u_0 u_1) + \order(\epsilon^2) = \epsilon 2u_0\vnabla u_1 + \order(\epsilon^2) ,\\
                \vnabla v\cdot \vnabla (u^2) &= \order(\epsilon^2),\\
                \nabla^2  u  &= \epsilon \nabla^2 u_1 \implies v\nabla^2  u=  \epsilon v_0 \nabla^2 u_1 +  \order(\epsilon^2).
              \end{align}
              So to $\order(\epsilon)$ our equations read
              \begin{align}
                \pd{u_1}{t} &= \nabla^2 u_1 + \gamma\left(F_u u_1 + F_v v_1\right),\\
                \pd{v_1}{t} &= D\nabla^2 v_1 + \gamma\left(G_u u_1 + G_v v_1\right) + v_0 \nabla^2 u_1,
              \end{align}
              where $F_u = \pd{F}{u}$ evaluated at $u=u_0$, etc.

              We seek solutions in the form $u_1 = \Re\left(A\e^{\i\v{k}\cdot  x + \lambda t}\right)$ and $v_1 = \Re\left(B\e^{\i\v{k}\cdot  x + \lambda t}\right)$. Our linearised equations then take the form
              \begin{equation}
                \lambda
                \begin{pmatrix}u\\v
                \end{pmatrix} = \t{J}
                \begin{pmatrix}u\\v
                \end{pmatrix},
              \end{equation}
              where
              \begin{equation}
                \t{J} = \mat{-k_s+ \gamma F_u }{\gamma F_v}{-k_sv_0  + \gamma G_u}{-Dk_s + \gamma G_v}
              \end{equation}
              and where $k_s = k_1^2 + k_2^2$.

              In order to have nontrivial $u_1,v_1$ we require $\det(\m{J}-\lambda \m{I})=0$ so
              \begin{equation}
                \lambda(i,j) = \frac{1}{2}\left[\tr\m{J} \pm \sqrt{\left(\tr\m{J}\right)^2 - 4\det\m{J}}\right]
              \end{equation}
              with
              \begin{align}
                \tr\m{J} &= \gamma(F_u + G_v)  - k_s(D+1),\\
                \det\m{J} & =D k_s^2 - \gamma k_s\left(D F_u + G_v - v_0 F_v\right)  + \gamma^2(F_uG_v- F_vG_u).
              \end{align}
              The boundary conditions require
              \begin{equation}
                k_s = \pi^2\left(\frac{n_1^2}{L_1^2}+\frac{n_2^2}{L_2^2}\right).
              \end{equation}
              First we require the $k_s=0$ mode be asymptotically stable (decay) so we require $\tr\m{J}<0$ and $\det\m{J}>0$ when $k_s=0$, i.e.
              \begin{equation}
                F_u + G_v<0,\mbox{ and }F_uG_v- F_vG_u>0.
              \end{equation}
              Next we want some $k_s>0$ such that some modes $(n_1,n_2)$ grow. Since $(F_u + G_v)<0$ the trace must be negative, so this can only occur if $\det\m{J}<0$. Since $F_uG_v- F_vG_u>0$ we require
              \begin{equation}
                D F_u + G_v - v_0 F_v >0,
              \end{equation}
              a necessary but not sufficient condition.

              Finally, since the quadratic component of $k_s$ is positive we can guarantee solutions if the quadratic $\det\m{J}$ crosses the real axis, thus the discriminant of the quadratic $\det\m{J}$ must be greater than $0$, i.e.
              \begin{equation}
                \left(D F_u + G_v - v_0 F_v\right)^2 - 4D(F_uG_v- F_vG_u) >0.
              \end{equation}
            }
          \item{
              If we set $u=0$ we get $v=0,1$. If we set $v=0$ we get $u=0,1$. So our equilibria are $(0,0)$, $(1,0)$ and $(0,1)$.
            }
          \item{
              We have
              \begin{equation}
                F_u = 1-2u_0 -v_0,\quad F_v = -u_0,\quad G_v = 1-2v_0+\gamma u_0,\quad G_u = \gamma v_0.
              \end{equation}

              Consider the $(0,0)$ equilibrium: We try the first condition
              \begin{equation}
                F_u + G_v = 2 >0
              \end{equation}
              so we cannot satisfy the first condition.

              Consider the $(1,0)$ equilibrium: The first condition reads
              \begin{equation}
                F_u + G_v =\gamma  <0,
              \end{equation}
              so we require that $\gamma$ be negative. The second condition is
              \begin{equation}
                F_u G_v - F_v G_u= -(1+\gamma) >0
              \end{equation}
              so this requires that $\gamma$ is less than $-1$. The third condition reads
              \begin{equation}
                D F_u + G_v =-D + (1+\gamma) >0 \implies \gamma > D-1.
              \end{equation}
              We see that it is not possible to satisfy $\gamma<-1$ and $\gamma> D-1$. So we cannot satisfy all the Turing conditions.

              Now consider the $(0,1)$ equilibrium:
              \begin{equation}
                F_u + G_v  =-1<0.
              \end{equation}
              The second condition requires
              \begin{equation}
                F_u G_v - F_v G_u  = 0 >0
              \end{equation}
              which cannot be satisfied.

              So in fact we cannot have any pattern formation from homogeneous equilibria for which one of the densities is zero.
            }
        \end{enumerate}

      }
      %

      %===========================================================================

      %
      \question{Patterns and polar coordinates (Exam section A, 2017)}
      {
        Consider the following \emph{non-dimensionalised} reaction--diffusion equation for the scalar densities $u$ and $v$
        \begin{subequations}
          \label{rds}
          \begin{align}
            \pd{u}{t} = \nabla^2 u + \gamma F(u,v),\\
            \pd{v}{t} = D \nabla^2 v + \gamma G(u,v),
          \end{align}
        \end{subequations}
        where $D>0$ and $\gamma>0$ are constants. We assume \emph{no-flux} boundary conditions on a Cartesian domain $[0,L_1]\times[0,L_2]$ with coordinates $(x_1,x_2)$. The Turing conditions for pattern formation are
        \begin{align}
          & F_u +G_v <0, \quad F_uG_v  -G_uF_v> 0,\nonumber \\
          & G_v + DF_u>0,\quad (G_v+ D F_u)^2 - 4 D (F_uG_v  -G_uF_v) >0, \label{turingconditions3}
        \end{align}
        where we have used the notation
        \begin{align*}
          &F_{u} = \pd{F}{u}\bigg\vert_{u=u_0,v=v_0},\quad F_{v} = \pd{F}{v}\bigg\vert_{u=u_0,v=v_0},\\
          &G_{u} = \pd{G}{u}\bigg\vert_{u=u_0,v=v_0},\quad G_{v} = \pd{G}{v}\bigg\vert_{u=u_0,v=v_0},
        \end{align*}
        with $(u_0,v_0)$ a homogeneous equilibrium of the system.
        \begin{enumerate}[label=(\alph*)]
          \item State the conditions for a homogeneous equilibrium of \cref{rds}.
          \item Define what is meant by a pattern in this context, give the explicit mathematical form for the pattern and describe what the conditions in \cref{turingconditions3} enforce.

          \item What would be the mathematical form of the patterns if the domain had instead been an annulus, covered with a polar coordinate system $(r,\theta)$? You do \emph{not} need to state the condition on the pattern numbers $k_s$ which results from the boundary conditions.

        \end{enumerate}
      }
      {
        \begin{enumerate}[label=(\alph*)]
          \item The homogeneous equilibrium satisfies
            \begin{equation}
              F(u,v) =0,\quad G(u,v) = 0.
            \end{equation}

          \item{
              This homogeneous equilibrium is unstable to small variations in density. This instability develops into an inhomogeneous pattern in the form often of spots or stripes (this being a Cartesian grid). The linearised version of the main equations has solutions in the form
              \begin{equation}
                \sum_{i=0}^{\infty}\sum_{j=0}^{\infty}\e^{\lambda(k_s^{ij}) t}\cos\left(\frac{i\pi}{L_1}x_1\right)\cos\left(\frac{j\pi}{L_2}x_2\right),
              \end{equation}
              where
              \begin{equation}
                k_s^{ij} = \frac{i^2\pi^2}{L_1^2}+\frac{j^2\pi^2}{L_2^2}.
              \end{equation}
              The function $\lambda(k_s)$ dictates if the mode decays exponentially $\lambda(k_s)<0,$ or grows, $\lambda(k_s)\geq 0$. The growing modes form the pattern.

              The first two conditions, $F_u +G_v <0$ and $F_uG_v  -G_uF_v> 0$, ensure that the $k_s^{00}=0$ mode, the homogeneous perturbation, always decays.

              The other two conditions, $G_v + DF_u>0$ and $(G_v+ D F_u)^2 - 4 D (F_uG_v  -G_uF_v)>0$, ensure there is some finite set of wave modes $k_s\in[k_s^{\min},k_s^{\max}]$ for which $\lambda(k_s^{ij})>0$ and hence a pattern can form.
            }
          \item{
              The patterns (solutions to the linearised system) take the form
              \begin{align}
                \sum_{n=0}^{\infty}\sum_{i=1}^{\infty}\left\{\left[A_{ni}J_n(\sqrt{k_{sn}^{i}}r)  + B_{ni}Y_n(\sqrt{k_{sn}^{i}}r) \right]\cos(n \theta) \right. \qquad \nonumber\\ + \left.\left[C_{ni}J_n(\sqrt{k_{sn}^{i}}r)  + E_{ni}Y_n(\sqrt{k_{sn}^{i}}r) \right]\sin(n\theta)\right\}.
              \end{align}
              Here $J_n$ and $Y_n$ are the Bessel functions of the first and second kind respectively, $A_{ni}, B_{ni}, C_{ni}$ and $E_{ni}$ are constants and $k_{sn}^{i}$ is the $i^{\rm th}$ pattern number associated with the $n^{\rm th}$ Fourier mode.
            }
        \end{enumerate}

      }
      %

      %===========================================================================

      %
      \question{More Turing analysis (Exam section B, 2017)}
      {
        Consider the following generic  reaction--diffusion system
        \begin{align*}
          \pd{u}{t} &= \nabla^2 u + \gamma F(u,v) + \nabla^2{uv},\\
          \pd{v}{t} &= D\nabla^2 v + \gamma G(u,v),
        \end{align*}
        where $D$ is the ratio of diffusion of $v$ and $u$ and $\gamma$ is a positive constant. We assume that nothing can leave or enter the system and that the domain $V$ is Cartesian.
        \begin{enumerate}[label=(\alph*)]
          \item{State the boundary conditions of this system.}
          \item{State the equations which would determine the homogeneous equilibria $(u_0,v_0)$ of this system.
            }
          \item{Perform the steps of the Turing analysis for this system to show that there will be a growing inhomogeneous pattern for the system when
              \begin{align*}
                &F_u + G_v <0, \quad F_u G_v - F_vG_u >0,\\
                &G_v(1+v_0) + D F_u - G_u u_0 >0,\\
                &\left[G_v(1+v_0) + D F_u - G_u u_0\right]^2 - 4 D(1+v_0)(F_u G_v - F_vG_u) \geq 0,
              \end{align*}
              where we have used the notation
              \begin{align*}
                &F_{u} = \pd{F}{u}\bigg\vert_{u=u_0,v=v_0},\quad F_{v} = \pd{F}{v}\bigg\vert_{u=u_0,v=v_0}\\
                &G_{u} = \pd{G}{u}\bigg\vert_{u=u_0,v=v_0},\quad G_{v} = \pd{G}{v}\bigg\vert_{u=u_0,v=v_0}.
              \end{align*}
            }
          \item{
              Now we set
              \begin{equation}
                F  = u\left(\frac{k_1 v}{k_2+v}-u\right),\quad G = (u-1)(v-1),
              \end{equation}
              where the constants $k_1$ and $k_2$ are positive. Find the three equilibria of the system and in each case state whether they can satisfy the inequalities
              \begin{align*}
                &F_u + G_v <0, \quad F_u G_v - F_vG_u >0,\\
                &G_v(1+v_0) + D F_u - G_u u_0 >0.
              \end{align*}
            }
        \end{enumerate}
      }
      {
        \begin{enumerate}[label=(\alph*)]
          \item The boundary conditions are
            $\vnabla u \cdot \uv{n} =0$ and $\vnabla v \cdot \uv{n}$, where $\uv{n}$ is the unit normal to $\partial{V}$.

          \item The homogeneous equilibria are determined by
            $F(u_0,v_0)=0$ and $G(u_0,v_0)=0$.

          \item{
              We linearise our system by Taylor expanding $u \approx u_0 + \epsilon u_1$ and $v \approx v_0 +\epsilon v_1$. Denoting $F_{u} = \pd{F}{u}\vert_{u=u_0,v=v_0}$ etc.\ we have, at order $\epsilon$,
              \begin{align*}
                \pd{u_1}{t} &= \nabla^2 u_1 + \gamma F_u u_1 + \gamma F_v v_1 + u_0\nabla^2v_1 + v_0\nabla^2u_1,\\
                \pd{v_1}{t} &= D\nabla^2 v_1 + \gamma G_u u_1 + \gamma G_v v_1.
              \end{align*}
              Assuming $u_1 =C_u\e^{\i \v{k}\cdot \v{x} + \lambda t}$ and $v_1 =C_v\e^{\i \v{k}\cdot \v{x} + \lambda t}$. We can write this as a matrix equation
              \begin{equation}
                \lambda \mvec{u_1}{v_1} = \m{J}\mvec{u_1}{v_1},
              \end{equation}
              where
              \begin{equation}
                \m{J} = \mat{-k_s(1+v_0) +\gamma F_u}{-k_s u_0+\gamma F_v}{\gamma G_u}{-Dk_s+ \gamma G_v}
              \end{equation}
              and $k_s = \v{k}\cdot\v{k}$. So solutions can only occur when $\det(\m{J}-\lambda \m{I})=0$. We know from the notes that solutions take the form
              \begin{equation}
                \lambda = \frac{1}{2}\left[\tr\m{J} \pm \sqrt{(\tr\m{J})^2 - 4\det\m{J}}\right],
              \end{equation}
              where
              \begin{align}
                \label{tures}
                \tr\m{J} & =  \gamma(F_u + G_v)- k_s(D+v_0+1),\\
                \nonumber \det\m{J}& =  D(1+v_0)k_s^2 - k_s\gamma\left(G_v(1+v_0) + D F_u - G_u u_0\right) + \gamma^2(F_u G_v - F_vG_u).
              \end{align}
              We first need the homogeneous perturbations $k_s=0$ to decay, thus we require $\tr <0$ and $\det>0$, so
              \begin{equation}
                F_u + G_v<0 \mbox{ and } F_uG_v - F_vG_u >0.
              \end{equation}
              This gives the first two conditions.

              Next we want there to be some $k_s>0$ for which $\lambda(k_s)>0$, so they grow. Since $F_u + G_v<0$ the trace must be negative ($v_0$ must be positive) so we can only have positive eigenvalues $\lambda$ when $\det<0$, i.e.
              \begin{equation}
                D(1+v_0)k_s^2 - k_s\gamma\left(G_v(1+v_0) + D F_u - G_u u_0\right) + \gamma^2(F_u G_v - F_vG_u) >0.
              \end{equation}
              Since $F_uG_v - F_vG_u>0$ and $D(1+v_0)>0$ this can only occur if
              \begin{equation}
                G_v(1+v_0) + D F_u - G_u u_0 >0,
              \end{equation}
              the third condition. But this is necessary not sufficient, finally, if we want real $k_s$ to solve $\det\m{J}=0$ we need the discriminant of the quadratic in $k_s$ to be positive, i.e.
              \begin{equation}
                \left(G_v(1+v_0) + D F_u - G_u u_0\right)^2 - 4 D(1+v_0)(F_u G_v - F_vG_u)\geq 0,
              \end{equation}
              the final condition.
            }
          \item{
              First $F=0$ when $u=0$, so $v=1$. We also solve $F=0$ by setting
              \begin{equation}
                \label{ucon}
                u = \frac{k_1v}{k_2+v}.
              \end{equation}
              So either $v=1$ or $u=1$ and \cref{ucon} gives the final two equilibria. The three equilibria $(u_0,v_0)$ are
              \begin{equation}
                (0,1),\quad \left(1,\frac{k_2}{k_1-1}\right) ,\quad \left(\frac{k_1}{1+k_2},1\right).
              \end{equation}
              We have
              \begin{align}
                F_u &= \frac{k_1 v_0}{k_2+ v_0}-2u_0, & F_v &= u_0\left(\frac{k_1}{k_2 +v_0}- \frac{k_1 v_0}{(k_2+v_0)^2}\right),\nonumber\\
                G_u &= v_0-1, & G_v & = u_0-1.
              \end{align}

              First the case $u_0=0$, $v_0=1$:
              \begin{equation}
                F_uG_v - F_vG_u =-\frac{k_1}{k_2+1} >0.
              \end{equation}
              This is a contradiction as $k_1,k_2>0$, so this equilibrium cannot satisfy \cref{tures}.

              Note for the next two cases, $u_0 = \frac{k_1 v_0}{k_2+ v_0}$, so
              \begin{equation}
                F_u = \frac{k_1 v_0}{k_2+ v_0}-2u_0 = -u_0.
              \end{equation}

              Next $(1,\frac{k_2}{k_1-1})$: The third condition gives
              \begin{equation}
                G_v(1+v_0) + D F_u - u_0G_u=  -D -(v_0-1) > 0.
              \end{equation}
              So at the very least we need $v_0<1$. Now the second condition
              \begin{equation}
                F_uG_v-F_vG_u = -F_vG_u >0.
              \end{equation}
              If $v_0<1$ then $G_u<0$ and we see that
              \begin{equation}
                F_v = \frac{u_0k_1}{k_2+v_0}\left(1- \frac{v_0}{k_2+v_0}\right).
              \end{equation}
              must be positive. So the second condition is satisfied if $v_0<1$. Finally the first condition
              \begin{equation}
                F_u + G_v= -u_0 <0.
              \end{equation}
              is satisfied. So if $v_0<1-D$, then the first three conditions will be satisfied.

              The final case is $(\frac{k_1}{k_2 +1},1)$. The second condition gives
              \begin{equation}
                F_uG_v - F_vG_u  = u_0(1-u_0)>0,
              \end{equation}
              so we require $u_0<1$.
              The third condition is
              \begin{equation}
                G_v(1+v_0) +DF_u - u_0G_u= (u_0-1)(1+v_0) -Du_0 > 0.
              \end{equation}
              But we needed $u_0<1$ so we cannot satisfy this condition.
            }
        \end{enumerate}

      }
      %

      %===========================================================================

      %
      \question{Boundary condition effects (Exam section A type)}
      {
        Consider the following \emph{non-dimensionalised} reaction--diffusion equation for the scalar densities $u(x,t)$ and $v(x,t)$,
        \begin{align*}
          \pd{u}{t} &= \nabla^2 u + \gamma F(u,v),\\
          \nonumber\pd{v}{t} &= D \nabla^2 v + \gamma G(u,v),
        \end{align*}
        where $D>0$ and $\gamma>0$ are constants. We assume \emph{no-flux} boundary conditions on a Cartesian domain $[0,L_1]\times[0,L_2]$ with coordinates $(x_1,x_2)$. The Turing conditions for pattern formation are
        \begin{align}
          & F_u +G_v <0, \quad F_uG_v  -G_uF_v> 0, \nonumber \\
          & G_v + DF_u>0,\quad (G_v+ D F_u)^2 - 4 D (F_uG_v  -G_uF_v) >0,\label{tcons}
        \end{align}
        where we have used the notation
        \begin{align*}
          &F_{u} = \pd{F}{u}\bigg\vert_{u=u_0,v=v_0},\quad F_{v} = \pd{F}{v}\bigg\vert_{u=u_0,v=v_0}\\
          &G_{u} = \pd{G}{u}\bigg\vert_{u=u_0,v=v_0},\quad G_{v} = \pd{G}{v}\bigg\vert_{u=u_0,v=v_0},
        \end{align*}
        with $(u_0,v_0)$ a homogeneous equilibrium of the system.
        \begin{enumerate}[label=(\alph*)]
          \item
            Define what is meant by a pattern in this context, give the explicit mathematical form for the pattern and describe what the conditions in \cref{tcons} enforce.

          \item Consider a system which is identical to the one considered in part (a), except with differing boundary conditions. Consider the following two (distinct) possibilities:
            \begin{enumerate}[label=(\roman*)]  \label{chemosys}
              \item{Homogeneous periodic boundary conditions along both directions.}
              \item{Fixed (homogeneous) values $u=u_0$ and $v=v_0$ on the $x_1=0,x_1=L$ boundaries (Dirichlet boundary conditions) and no-flux boundary conditions on the $x_2=0,x_2=L$ boundaries  (for both $u$ and $v$).}
            \end{enumerate}
            Argue in each case if the Turing conditions,
            \begin{equation}
              F_u +G_v <0, \quad F_uG_v  -G_uF_v> 0,
            \end{equation}
            are necessary for pattern formation as defined in part (a).

        \end{enumerate}
      }
      {
        \begin{enumerate}[label=(\alph*)]
          \item
            The homogeneous equilibrium defined by $F(u,v) =0,\,G(u,v) = 0$ is unstable to small variations in density. This instability develops into an inhomogeneous pattern in the form of spots or stripes (this being a Cartesian grid). The linearised version of the main equations has solutions of the form
            \begin{equation}
              \sum_{i=0}^{n}\sum_{j=0}^{m}\e^{\lambda(k_s^{ij}) t}\cos\left(\frac{n_i\pi}{L_1}x_1\right)\cos\left(\frac{m_j\pi}{L_2}x_2\right),
            \end{equation}
            where
            \begin{equation}
              k_s^{ij} = \frac{n_i^2\pi^2}{L_1^2}+\frac{m_j^2\pi^2}{L_2^2}.
            \end{equation}
            The function $\lambda(k_s)$ dictates if the mode decays exponentially, $\lambda(k_s)<0$, or grows, $\lambda(k_s)\geq 0$. The growing modes form the pattern.

            The first two conditions, $F_u +G_v <0$ and $F_uG_v  -G_uF_v> 0$, ensure that the $k_s^{00}=0$ mode, the homogeneous perturbation, always decays. The other two conditions $G_v + DF_u>0$ and $(G_v+ D F_u)^2 - 4 D (F_uG_v  -G_uF_v)>0$, ensure there is some finite set of wave modes, $k_s\in[k_s^{\min},k_s^{\max}]$, for which $\lambda(k_s^{ij})>0$, and hence a pattern can form.

          \item
            In the case of periodic boundary conditions both the $\sin(k_{1,2}x)$ and $\cos(k_{1,2}x)$ spatial functions are permissible solutions. This means the $n=m=0$ ($k_s=0$) homogeneous case is permissible and the two stated conditions are necessary. In the second case the fixed-value boundary conditions at $x_1=0,x_1=L$ means we can have no $n_1=0$ mode, thus no $k_s=0$ pattern.  This means the two conditions are unnecessary.

        \end{enumerate}

      }
      %

      %===========================================================================

      %
      \question{Enzyme-driven Turing patterns (Exam section A type)}
      {
        Consider the following reaction--diffusion model of hair growth in the \emph{Acetabularia} sea plant. The hair growth hormone $u$ and calcium concentration $v$ which interact to stimulate tip growth is represented by the following set of partial differential equations:
        \begin{subequations}
          \begin{align}
            \pd{u}{t} &= \nabla^2 u + \gamma(a - u + u^2v),\\
            \pd{v}{t} &= D\nabla^2 v + \gamma(b- u^2 v).
          \end{align}
        \end{subequations}
        Here $D,\gamma,a$ and $b$ are positive constants and we impose no-flux boundary conditions. It is assumed spontaneous tip growth occurs due to a Turing instability of the system. We consider this system on an annular domain, with no-flux boundary conditions imposed.
        \begin{enumerate}[label=(\alph*)]
          \item{The homogeneous equilibria of the system are:
              \begin{equation}
                u_0 =a+b,\quad v_0 = u_0-a.
              \end{equation}
              The Turing conditions for pattern formation are
              \begin{align}
                & F_u +G_v <0, \quad F_uG_v  -G_uF_v> 0,\nonumber \\
                & G_v + DF_u>0, \label{turingconditions4}
              \end{align}
              and
              \begin{align*}
                (G_v+ D F_u)^2 - 4 D (F_uG_v  -G_uF_v) >0,
              \end{align*}
              where we have used the notation
              \begin{align*}
                F_{u} &= \pd{F}{u}\bigg\vert_{u=u_0,v=v_0},& F_{v} &= \pd{F}{v}\bigg\vert_{u=u_0,v=v_0},\\
                G_{u} &= \pd{G}{u}\bigg\vert_{u=u_0,v=v_0},& G_{v} &= \pd{G}{v}\bigg\vert_{u=u_0,v=v_0},
              \end{align*}
              with $(u_0,v_0)$ a homogeneous equilibrium of the system. In this case we have
              \begin{align*}
                F_u &=2u_0v_0 -1 = 1- \frac{2a}{u_0},& F_v &= u_0^2,\\
                G_u&=-2u_0v_0 = -\frac{2(u_0-a)}{u_0},& G_v &= -u_0^2.
              \end{align*}

              Find conditions on $a$ such that the first three Turing conditions, \cref{turingconditions4}, are satisfied.
            }
          \item{In addition, the fourth Turing condition can be shown to yield the following pair of constraints:
              \begin{equation}
                a< \frac{u_0}{2}\left(1-\frac{2u_0}{\sqrt{D}} -\frac{u_0^2}{D}\right) \quad \text{or} \quad a> \frac{u_0}{2}\left(1+\frac{2u_0}{\sqrt{D}} -\frac{u_0^2}{D}\right).
              \end{equation}

              Assume $D>1$. Argue that if $a$ satisfies the following two inequalities,
              \begin{equation}
                a<  \frac{u_0}{2}\left(1-\frac{2u_0}{\sqrt{D}} -\frac{u_0^2}{D}\right) \quad \text{and} \quad a>  \frac{u_0(1-u_0^2)}{2},
              \end{equation}
              then all four Turing  conditions are satisfied.
            }
        \end{enumerate}
      }
      {
        \begin{enumerate}[label=(\alph*)]
          \item{
              The first condition gives
              \begin{equation}
                \label{acon}
                1-\frac{2a}{u_0} - u_0^2<0 \implies a> \frac{u_0(1-u_0^2)}{2}.
              \end{equation}

              The second condition is trivially satisfied,
              \begin{equation}
                F_uG_v -F_vG_u = u_0^2>0,
              \end{equation}
              so no condition on $a$.

              The third condition gives
              \begin{equation}
                \label{acon3}
                D(1-\frac{2a}{u_0}) -u_0^2\implies a<\frac{u_0}{2}\left(1-\frac{u_0^2}{D}\right).
              \end{equation}
            }
          \item{
              As
              \begin{equation}
                \frac{u_0}{2}\left(1-\frac{2u_0}{\sqrt{D}} -\frac{u_0^2}{D}\right)<  \frac{u_0}{2}\left(1-\frac{u_0^2}{D}\right),
              \end{equation}
              we see that \cref{acon3} is satisfied if
              \begin{equation}
                a<  \frac{u_0}{2}\left(1-\frac{2u_0}{\sqrt{D}} -\frac{u_0^2}{D}\right).
              \end{equation}
              Next,
              \begin{equation}
                \frac{u_0}{2}\left(1+\frac{2u_0}{\sqrt{D}} -\frac{u_0^2}{D}\right) > \frac{u_0}{2}\left(1-u_0^2\right),
              \end{equation}
              so satisfying the other condition 4 inequality,
              \begin{equation}
                a> \frac{u_0}{2}\left(1+\frac{2u_0}{\sqrt{D}} -\frac{u_0^2}{D}\right),
              \end{equation}
              ensures \cref{acon} is satisfied. But we must check \emph{each} condition 4 inequality satisfies \emph{both} conditions 1 and 3.

              We also note that
              \begin{equation}
                \frac{u_0}{2}\left(1+\frac{2u_0}{\sqrt{D}} -\frac{u_0^2}{D}\right) >   \frac{u_0}{2}\left(1-\frac{u_0^2}{D}\right)> \frac{u_0}{2}\left(1-\frac{2u_0}{\sqrt{D}} -\frac{u_0^2}{D}\right).
              \end{equation}
              So it is not possible to satisfy
              \begin{equation}
                a>   \frac{u_0}{2}\left(1+\frac{2u_0}{\sqrt{D}} -\frac{u_0^2}{D}\right).
              \end{equation}
              and
              \begin{equation}
                a<\frac{u_0}{2}\left(1-\frac{u_0^2}{D}\right).
              \end{equation}
              That is to say one of the possible cases for condition 4 holding would fail to also satisfy condition 3. We are left with only one possibly viable inequality to satisfy condition 4,
              \begin{equation}
                a< \frac{u_0}{2}\left(1-\frac{2u_0}{\sqrt{D}} -\frac{u_0^2}{D}\right),
              \end{equation}
              which we know from above satisfies condition 3. We need to check it can satisfy condition 1.

              Indeed we can see it is possible that
              \begin{equation}
                \frac{u_0}{2}\left(1-u_0^2\right) <  \frac{u_0}{2}\left(1-\frac{2u_0}{\sqrt{D}} -\frac{u_0^2}{D}\right).
              \end{equation}
              in the limit $D\to \infty$.
            }
        \end{enumerate}

      }
  \end{enumerate}
  \end{document}
